% ═══════════════════════════════════════════════════════════════════
%  Chapitre 2 — Analyse et conception
% ═══════════════════════════════════════════════════════════════════
\chapter{Analyse et conception}

\section*{Introduction}
\addcontentsline{toc}{section}{Introduction}

Ce chapitre présente le résultat du travail d'analyse et de conception. Il
commence par la spécification des besoins --- fonctionnels puis non
fonctionnels --- et par l'identification des acteurs du système. Il expose
ensuite la modélisation UML retenue : les diagrammes de cas d'utilisation,
leur description textuelle, les diagrammes de séquence des processus les plus
délicats, et enfin le diagramme de classes qui fixe la structure des données.

Une précision liminaire sur la méthode : les modèles présentés ici ne sont pas
des artefacts produits en amont puis abandonnés. Ils ont été révisés à chaque
sprint, et le diagramme de classes en particulier reflète l'état final d'un
schéma qui a connu vingt migrations successives. Les choix qui l'ont fait
évoluer sont explicités en fin de chapitre, car ils constituent l'essentiel de
l'apport de conception de ce projet.

% ─────────────────────────────────────────────────────────────────
\section{Spécifications des besoins}

\subsection{Besoins fonctionnels}

Les besoins fonctionnels décrivent ce que le système doit permettre de faire.
Ils ont été regroupés en huit domaines.

\paragraph{Gestion des comptes et de l'authentification.}
\begin{itemize}[leftmargin=1.2cm,nosep]
  \item Inscription d'une cliente avec adresse e-mail, mot de passe, nom,
  prénom et téléphone.
  \item Confirmation de l'adresse e-mail par un code à usage unique.
  \item Connexion, déconnexion et maintien de session sans reconnexion
  quotidienne.
  \item Réinitialisation du mot de passe oublié, par code envoyé par e-mail.
  \item Consultation et modification du profil, changement de mot de passe.
  \item Suppression du compte à la demande de la cliente.
\end{itemize}

\paragraph{Catalogue de prestations.}
\begin{itemize}[leftmargin=1.2cm,nosep]
  \item Consultation publique des prestations actives, par catégorie.
  \item Consultation du détail d'une prestation : description, tarif, photo.
  \item Création, modification, activation et désactivation des prestations et
  de leurs catégories par la gérante.
\end{itemize}

\paragraph{Catalogue de produits.}
\begin{itemize}[leftmargin=1.2cm,nosep]
  \item Consultation publique des produits actifs, par catégorie.
  \item Affichage de la disponibilité en stock.
  \item Gestion complète du catalogue et des stocks par la gérante.
\end{itemize}

\paragraph{Gestion des rendez-vous.}
\begin{itemize}[leftmargin=1.2cm,nosep]
  \item Consultation des créneaux disponibles pour une prestation à une date
  donnée.
  \item Réservation d'un créneau par la cliente.
  \item Consultation de ses propres rendez-vous, à venir et passés.
  \item Annulation par la cliente.
  \item Réservation \emph{pour le compte} d'une cliente par le personnel, pour
  les rendez-vous pris par téléphone ou au comptoir.
  \item Confirmation, clôture, annulation et report d'un rendez-vous par le
  personnel.
\end{itemize}

\paragraph{Gestion des commandes.}
\begin{itemize}[leftmargin=1.2cm,nosep]
  \item Panier et passage de commande, avec retrait à l'institut ou livraison à
  domicile.
  \item Paiement à la remise ou à la livraison.
  \item Consultation de ses commandes et de leur avancement par la cliente.
  \item Annulation par la cliente tant que la commande n'a pas été remise.
  \item Traitement des commandes par le personnel selon un cycle de vie
  contrôlé.
\end{itemize}

\paragraph{Programme de fidélité.}
\begin{itemize}[leftmargin=1.2cm,nosep]
  \item Gain automatique de points sur les prestations réalisées et les
  commandes remises.
  \item Consultation du solde et de l'historique des mouvements.
  \item Catalogue de récompenses échangeables contre des points.
  \item Paliers de visites offrant automatiquement une récompense, avec option
  de récurrence.
  \item Émission d'un bon de récompense doté d'un code lisible, honoré au
  comptoir.
  \item Ajout manuel de points par le personnel, avec motif obligatoire et
  journal consultable par la gérante.
\end{itemize}

\paragraph{Paramétrage du centre.}
\begin{itemize}[leftmargin=1.2cm,nosep]
  \item Définition des horaires d'ouverture, à raison de plusieurs plages par
  jour.
  \item Déclaration de fermetures exceptionnelles (congés, jours fériés).
  \item Réglage du nombre de places réservables simultanément et de l'écart
  entre créneaux, sans redéploiement.
\end{itemize}

\paragraph{Pilotage et exports.}
\begin{itemize}[leftmargin=1.2cm,nosep]
  \item Tableau de bord : commandes en attente, rendez-vous du jour, alertes de
  stock faible.
  \item Gestion des comptes du personnel et attribution des rôles.
  \item Export au format Excel des clientes, des commandes et des rendez-vous,
  réservé à la gérante.
\end{itemize}

\subsection{Besoins non fonctionnels}

Les besoins non fonctionnels portent sur les qualités attendues de la solution
plutôt que sur ses fonctions. Sept ont été retenus, chacun assorti du moyen
concret par lequel il est satisfait --- un besoin non fonctionnel qui ne se
traduit par aucune décision d'implémentation n'est qu'une intention.

\begin{table}[H]
\centering
\caption{Besoins non fonctionnels et moyens de satisfaction}
\label{tab:besoins-non-fonctionnels}
\begin{tabularx}{\textwidth}{L{3.1cm} X}
\toprule
\entete{Besoin} & \entete{Moyen retenu} \\
\midrule
\textbf{Exactitude} &
Aucun créneau ni aucune unité de stock ne peut être attribué deux fois. Les
réservations concurrentes sont sérialisées par un verrou consultatif
PostgreSQL ; le stock et le solde de points sont modifiés par écritures
atomiques conditionnelles. \\

\textbf{Sécurité} &
Mots de passe hachés avec Argon2 ; jetons d'accès JWT de courte durée ; jetons
de rafraîchissement stockés hachés, avec rotation et détection de réutilisation ;
autorisation par rôle vérifiée côté serveur sur chaque route ; en-têtes de
sécurité HTTP ; validation stricte de toute entrée ; limitation du débit par
adresse IP. \\

\textbf{Fiabilité} &
Aucune erreur technique n'atteint l'utilisatrice : un filtre d'exceptions
global traduit les erreurs de base de données en messages compréhensibles.
Les états des rendez-vous et des commandes sont régis par des machines à états
explicites qui interdisent toute transition incohérente. \\

\textbf{Performance} &
Toutes les listes susceptibles de croître sont paginées côté serveur ; le
tableau de bord ne demande que ce qu'il affiche ; les images du catalogue sont
servies avec un cache d'un an et décodées à la taille réellement affichée. \\

\textbf{Ergonomie} &
Charte graphique unifiée entre les trois applications, modes clair et sombre,
messages en langue naturelle, aucun terme technique exposé. \\

\textbf{Accessibilité linguistique} &
Interface mobile disponible en français, en arabe et en anglais ; les réponses
de l'API suivent l'en-tête \texttt{Accept-Language} de la requête et les
e-mails la langue choisie par la cliente. \\

\textbf{Conformité} &
Droit à l'effacement réellement exerçable, par anonymisation du compte plutôt
que par suppression en cascade, afin de préserver l'historique comptable de
l'institut (loi 09-08, CNDP). \\
\bottomrule
\end{tabularx}
\end{table}

\subsection{Acteurs du système}

Cinq acteurs interagissent avec le système : quatre acteurs humains --- la
visiteuse non authentifiée et les trois rôles applicatifs --- et un acteur
secondaire non humain.

\begin{itemize}[leftmargin=1.2cm]
  \item \textbf{La visiteuse} --- acteur non authentifié. Elle consulte les
  catalogues publics et peut créer un compte. Elle ne peut ni réserver ni
  commander.

  \item \textbf{La cliente} (\texttt{CLIENT}) --- acteur principal de
  l'application mobile. Elle réserve, commande, suit son compte de fidélité et
  gère son profil. Elle n'a aucun accès au back-office : une tentative de
  connexion y est refusée à l'entrée.

  \item \textbf{Le personnel} (\texttt{STAFF}) --- acteur principal du
  back-office. Il traite les commandes et les rendez-vous, réserve pour le
  compte d'une cliente, honore les bons de récompense et crédite manuellement
  des points.

  \item \textbf{La gérante} (\texttt{ADMIN}) --- elle dispose de tous les
  droits du personnel, auxquels s'ajoutent la maîtrise du catalogue, des
  horaires, des réglages du centre, du programme de fidélité, des comptes
  utilisateurs, du journal des ajouts manuels de points et des exports Excel.

  \item \textbf{Le service de messagerie} --- acteur secondaire non humain. Il
  achemine les codes de vérification d'adresse et de réinitialisation de mot de
  passe.
\end{itemize}

\subsection{Langage de modélisation}

Le langage UML (\textit{Unified Modeling Language}) a été retenu pour
modéliser le système. C'est un langage graphique normalisé, indépendant de tout
processus de développement et de tout langage de programmation, ce qui permet
de l'intégrer sans friction à une démarche agile : les diagrammes sont révisés
à chaque sprint plutôt que figés au départ.

Trois types de diagrammes ont été employés, chacun pour ce qu'il montre le
mieux :

\begin{itemize}[leftmargin=1.2cm]
  \item les \textbf{diagrammes de cas d'utilisation}, pour délimiter le
  périmètre fonctionnel et répartir les droits entre acteurs ;
  \item les \textbf{diagrammes de séquence}, pour décrire l'ordre des échanges
  dans les processus où cet ordre est précisément ce qui fait la correction du
  système ;
  \item le \textbf{diagramme de classes}, pour fixer la structure des données
  persistantes et leurs relations.
\end{itemize}

% ─────────────────────────────────────────────────────────────────
\section{Diagrammes de cas d'utilisation}

\subsection{Diagramme de cas d'utilisation général}

La figure ci-dessous présente la vue d'ensemble des cas d'utilisation du
système, tous acteurs confondus.

\figProjet{cu-general.png}{Diagramme de cas d'utilisation général de la plate-forme}{cu-general}

\subsection{Cas d'utilisation : réservation d'une prestation}

Le domaine de la réservation est le plus riche du système. Il fait intervenir
la cliente et le personnel sur des cas partiellement communs : le personnel
peut réserver pour le compte d'une cliente, ce qui est le même cas
d'utilisation exercé par un autre acteur, mais il dispose en outre du report et
du changement de statut.

\figProjet{cu-reservation.png}{Diagramme de cas d'utilisation --- gestion des rendez-vous}{cu-reservation}

\subsection{Cas d'utilisation : commande de produits}

\figProjet{cu-commande.png}{Diagramme de cas d'utilisation --- commande de produits}{cu-commande}

\subsection{Cas d'utilisation : programme de fidélité}

Le programme de fidélité comporte deux circuits de récompense bien
distincts, qui convergent vers un même objet : le bon de récompense. Le premier
est automatique : un palier de visites franchi débloque une récompense que la
cliente réclame. Le second est volontaire : la cliente échange ses points
contre une récompense du catalogue. Dans les deux cas, un bon est émis, et
c'est ce bon que le personnel honore au comptoir.

\figProjet{cu-fidelite.png}{Diagramme de cas d'utilisation --- programme de fidélité}{cu-fidelite}

\subsection{Cas d'utilisation : administration}

\figProjet{cu-backoffice.png}{Diagramme de cas d'utilisation --- back-office}{cu-backoffice}

% ─────────────────────────────────────────────────────────────────
\section{Description des cas d'utilisation}

Les cas d'utilisation les plus structurants sont décrits ci-après sous forme
textuelle. Les scénarios alternatifs y occupent une place volontairement
importante : ce sont eux qui déterminent la robustesse du système, et plusieurs
d'entre eux n'ont été identifiés qu'à l'usage.

\begin{table}[H]
\centering
\caption{Description du cas d'utilisation « Réserver une prestation »}
\label{tab:cu-reserver}
\begin{tabularx}{\textwidth}{L{3.1cm} X}
\toprule
\entete{Cas d'utilisation} & Réserver une prestation \\
\midrule
\entete{Acteur principal} & La cliente \\
\entete{Objectif} & Obtenir un rendez-vous pour une prestation, à une date et
une heure disponibles \\
\entete{Préconditions} & La cliente est authentifiée ; la prestation est active
au catalogue \\
\entete{Scénario nominal} &
\begin{enumerate}[leftmargin=*,nosep]
  \item La cliente sélectionne une prestation et une date.
  \item Le système calcule la grille des créneaux du jour à partir des plages
  d'ouverture, de l'écart réglé par le centre et des fermetures
  exceptionnelles.
  \item Le système marque comme indisponible tout créneau dont la capacité est
  atteinte, ainsi que tout créneau déjà passé.
  \item La cliente choisit un créneau libre et valide.
  \item Le système prend le verrou de réservation, contrôle à nouveau la
  disponibilité, enregistre le rendez-vous au statut \textit{en attente} et fige
  le tarif de la prestation.
  \item Le système affiche la confirmation.
\end{enumerate} \\
\entete{Alternatives} &
\textbf{2a.} Le jour est fermé (aucune plage, ou fermeture exceptionnelle) : le
système l'indique et ne propose aucun créneau.\newline
\textbf{4a.} Le créneau a été pris entre l'affichage et la validation : le
système refuse et invite à en choisir un autre.\newline
\textbf{5a.} La date demandée est passée : le système refuse la réservation. \\
\entete{Postconditions} & Le rendez-vous existe au statut \textit{en attente} ;
la capacité du créneau est diminuée d'une unité \\
\bottomrule
\end{tabularx}
\end{table}

\begin{table}[H]
\centering
\caption{Description du cas d'utilisation « Passer une commande »}
\label{tab:cu-commander}
\begin{tabularx}{\textwidth}{L{3.1cm} X}
\toprule
\entete{Cas d'utilisation} & Passer une commande de produits \\
\midrule
\entete{Acteur principal} & La cliente \\
\entete{Objectif} & Commander des produits, avec retrait à l'institut ou
livraison à domicile \\
\entete{Préconditions} & La cliente est authentifiée ; son panier n'est pas
vide \\
\entete{Scénario nominal} &
\begin{enumerate}[leftmargin=*,nosep]
  \item La cliente ajoute des produits à son panier et ouvre la validation.
  \item Elle choisit le mode de remise ; en cas de livraison, elle saisit une
  adresse.
  \item Le système fusionne les lignes en double, contrôle que chaque produit
  est actif et que le stock est suffisant.
  \item Le système fige le prix unitaire de chaque ligne, calcule le total et
  enregistre la commande au statut \textit{en attente}.
  \item La cliente reçoit le récapitulatif de sa commande.
\end{enumerate} \\
\entete{Alternatives} &
\textbf{2a.} Livraison choisie sans adresse : le système refuse.\newline
\textbf{3a.} Un produit a été retiré du catalogue ou son stock est
insuffisant : le système refuse la commande en désignant le produit en
cause. \\
\entete{Postconditions} & La commande existe au statut \textit{en attente} ;
\emph{le stock n'est pas encore décrémenté} --- il ne l'est qu'à la
confirmation par le personnel \\
\bottomrule
\end{tabularx}
\end{table}

Ce dernier point mérite d'être souligné, car il résulte d'un arbitrage
explicite. Décrémenter le stock dès la commande aurait permis de le réserver,
mais aurait ouvert la porte au blocage du stock par des commandes jamais
honorées. Le paiement s'effectuant à la remise, l'institut confirme chaque
commande par un appel : c'est ce moment, et non la commande elle-même, qui
engage réellement le produit.

\begin{table}[H]
\centering
\caption{Description du cas d'utilisation « Échanger des points contre une récompense »}
\label{tab:cu-echanger}
\begin{tabularx}{\textwidth}{L{3.1cm} X}
\toprule
\entete{Cas d'utilisation} & Échanger des points contre une récompense \\
\midrule
\entete{Acteur principal} & La cliente \\
\entete{Objectif} & Convertir un solde de points en récompense utilisable à
l'institut \\
\entete{Préconditions} & La cliente est authentifiée ; la récompense est active
au catalogue \\
\entete{Scénario nominal} &
\begin{enumerate}[leftmargin=*,nosep]
  \item La cliente consulte le catalogue des récompenses et son solde.
  \item Elle demande l'échange d'une récompense.
  \item Le système débite le solde de façon atomique, conditionnellement à sa
  suffisance.
  \item Le système enregistre le mouvement et émet un bon de récompense muni
  d'un code lisible.
  \item La cliente présente ce code au comptoir lors de sa prochaine visite.
\end{enumerate} \\
\entete{Alternatives} &
\textbf{3a.} Le solde est devenu insuffisant entre l'affichage et la demande :
le débit conditionnel n'affecte aucune ligne, l'échange est refusé et rien
n'est émis. \\
\entete{Postconditions} & Le solde est diminué du coût de la récompense ; un
bon dû figure sur la liste du comptoir \\
\bottomrule
\end{tabularx}
\end{table}

\begin{table}[H]
\centering
\caption{Description du cas d'utilisation « Traiter une commande »}
\label{tab:cu-traiter}
\begin{tabularx}{\textwidth}{L{3.1cm} X}
\toprule
\entete{Cas d'utilisation} & Traiter une commande \\
\midrule
\entete{Acteur principal} & Le personnel \\
\entete{Objectif} & Faire progresser une commande jusqu'à sa remise à la
cliente \\
\entete{Préconditions} & Le membre du personnel est authentifié au back-office
avec le rôle \texttt{STAFF} ou \texttt{ADMIN} \\
\entete{Scénario nominal} &
\begin{enumerate}[leftmargin=*,nosep]
  \item Le personnel ouvre la liste des commandes en attente.
  \item Il confirme la commande après appel à la cliente : le système
  décrémente le stock de chaque ligne de façon atomique.
  \item Il marque la commande \textit{prête} lorsqu'elle est préparée.
  \item Il la marque \textit{terminée} à la remise : le système crédite les
  points de fidélité correspondants et met à jour le compteur de visites.
\end{enumerate} \\
\entete{Alternatives} &
\textbf{2a.} Le stock est devenu insuffisant : la confirmation échoue et la
commande reste en attente.\newline
\textbf{$\ast$a.} Une transition non autorisée est demandée (par exemple
\textit{terminée} depuis \textit{en attente}) : le système la refuse.\newline
\textbf{$\ast$b.} La commande est annulée après confirmation : le stock
décrémenté est restitué. \\
\entete{Postconditions} & La commande est au statut \textit{terminée} ; les
points sont crédités exactement une fois \\
\bottomrule
\end{tabularx}
\end{table}

% ─────────────────────────────────────────────────────────────────
\section{Diagrammes de séquence}

Cinq processus ont été modélisés sous forme de diagrammes de séquence. Ils ont
été choisis non pour leur fréquence d'usage, mais parce que dans chacun d'eux
\emph{l'ordre des opérations détermine la correction du résultat} --- ce qu'un
diagramme de cas d'utilisation ne peut pas montrer.

\subsection{Inscription et vérification de l'adresse e-mail}

\figProjet{seq-inscription.png}{Diagramme de séquence --- inscription et vérification de l'adresse}{seq-inscription}

L'inscription crée le compte, son compte de fidélité associé, et renvoie
immédiatement les jetons de session. Ce dernier point n'allait pas de soi : une
première version renvoyait un simple accusé de création, et l'application
enchaînait sur un appel de connexion. Entre les deux appels, la protection
anti-force brute pouvait refuser le second, laissant la cliente devant un
message d'échec alors même que son compte venait d'être créé. L'inscription
renvoie donc désormais exactement ce que renvoie la connexion.

Un code à six chiffres est envoyé par e-mail dans la langue choisie par la
cliente. Seule son empreinte est conservée en base : une fuite de la base ne
livrerait aucun code exploitable. Le nombre de tentatives est plafonné --- six
chiffres représentent un million de possibilités, ce qui se parcourt en
quelques minutes sans ce garde-fou. L'émission d'un nouveau code efface le
précédent du même usage, ce qui garantit qu'au plus une ligne active existe par
cliente et par usage.

\subsection{Réservation d'un créneau}

C'est le processus le plus délicat du système, et le diagramme ci-dessous en
montre la raison.

\figProjet{seq-reservation.png}{Diagramme de séquence --- réservation d'un créneau}{seq-reservation}

La difficulté tient à ce que la contrainte de capacité porte sur un
\emph{ensemble} de lignes --- le nombre de rendez-vous qui chevauchent le
créneau --- et non sur une ligne particulière. Aucune écriture conditionnelle
ne peut donc la garantir. La séquence « compter les rendez-vous du créneau,
puis insérer si la capacité le permet » est correcte pour une cliente seule, et
fausse dès que deux réservations se croisent : les deux comptent avant que
l'une ou l'autre n'ait inséré, les deux trouvent de la place, et deux cabines
en accueillent trois.

La solution retenue consiste à ouvrir une transaction et à y prendre un verrou
consultatif PostgreSQL, tenu jusqu'à la fin de la transaction. La seconde
réservation attend que la première soit écrite avant de compter. Ce verrou se
libère automatiquement au \textit{commit} comme au \textit{rollback}, ce qui
écarte tout risque de blocage résiduel.

À l'intérieur de cette transaction, la vérification enchaîne quatre contrôles :
la prestation existe et est active, la date n'est pas passée, le jour n'est ni
fermé ni exceptionnellement clos, et la capacité du créneau n'est pas atteinte.
Le tarif est enfin figé sur le rendez-vous, afin qu'une révision ultérieure des
prix ne modifie pas les points gagnés sur un rendez-vous déjà réservé.

\subsection{Traitement d'une commande et gestion du stock}

\figProjet{seq-commande.png}{Diagramme de séquence --- commande et décrément du stock}{seq-commande}

Le décrément du stock à la confirmation obéit à la même exigence d'atomicité
que la réservation, mais il se traite plus simplement : la contrainte porte
cette fois sur une ligne unique --- le produit --- et une mise à jour
conditionnelle suffit. La condition « décrémenter de $n$ à condition que le
stock soit au moins égal à $n$ » est évaluée par la base au moment de
l'écriture. Si elle n'est pas satisfaite, aucune ligne n'est affectée et la
transaction entière est annulée : la commande reste en attente et le stock
demeure cohérent.

\subsection{Clôture d'un rendez-vous et crédit de fidélité}

\figProjet{seq-fidelite.png}{Diagramme de séquence --- clôture d'un rendez-vous, crédit de points et palier}{seq-fidelite}

La clôture d'un rendez-vous déclenche une chaîne d'opérations qui doivent
toutes réussir ou toutes échouer : passage au statut \textit{terminé}, calcul
des points sur le tarif figé, incrément du solde et du compteur de visites,
écriture du mouvement d'historique, puis évaluation des paliers de visites. Le
tout s'exécute dans une seule transaction.

L'idempotence est assurée à deux niveaux. Le mouvement de fidélité porte une
référence unique vers le rendez-vous : une clôture rejouée ne peut pas créditer
deux fois. Les déblocages de paliers portent une contrainte d'unicité sur le
triplet (compte, palier, cycle), où le cycle indique combien de fois le seuil a
été franchi --- ce qui permet à un palier d'être récurrent tout en gardant
chaque déblocage unique.

\subsection{Rotation des jetons de session}

\figProjet{seq-refresh.png}{Diagramme de séquence --- rafraîchissement et rotation des jetons}{seq-refresh}

Le jeton d'accès est volontairement de courte durée. Pour éviter de demander à
la cliente de se reconnecter plusieurs fois par jour, un jeton de
rafraîchissement de longue durée permet d'en obtenir un nouveau.

Ce mécanisme est protégé par une rotation : chaque rafraîchissement révoque le
jeton présenté et en émet un nouveau, rattaché à la même famille. Si un jeton
déjà utilisé est présenté une seconde fois, c'est le signe qu'il a été
dérobé --- soit par l'attaquant, soit par la cliente légitime, sans qu'il soit
possible de savoir lequel. Le système révoque alors toute la famille : les deux
parties sont déconnectées, et l'accès frauduleux prend fin.

% ─────────────────────────────────────────────────────────────────
\section{Diagramme de classes}

\subsection{Vue d'ensemble}

Le diagramme de classes ci-dessous présente la structure des données
persistantes. Il comporte dix-neuf entités, réparties en six domaines :
utilisateurs et sessions, catalogues, rendez-vous, commandes, fidélité, et
paramétrage du centre.

\figProjet[0.98]{diagramme-classes.png}{Diagramme de classes de la plate-forme}{diagramme-classes}

\subsection{Choix structurants du modèle}

Cinq décisions de modélisation méritent d'être explicitées, car elles ne
découlent pas mécaniquement des besoins et constituent l'essentiel de l'apport
de conception.

\paragraph{Le figement des montants.} Un rendez-vous conserve le tarif de la
prestation \emph{tel qu'il était à la réservation}, et chaque ligne de commande
conserve le prix unitaire du produit \emph{tel qu'il était à la commande}.
Sans cela, une révision des tarifs réécrirait rétroactivement l'historique :
une commande passée en juillet afficherait le prix d'août, et les points gagnés
sur un rendez-vous changeraient après coup. Le même principe s'applique aux
récompenses de palier, dont la référence est figée au déblocage, et aux points
dépensés, figés sur le bon émis.

\paragraph{L'anonymisation plutôt que la suppression.} Une cliente qui exerce
son droit à l'effacement voit son compte vidé de tout ce qui l'identifie ---
adresse e-mail, nom, téléphone, mot de passe --- et marqué comme supprimé, mais
sa ligne n'est pas détruite. Ses commandes et ses rendez-vous portent en effet
le chiffre d'affaires de l'institut, qui doit survivre à son départ. Les
relations le garantissent structurellement : les rendez-vous et les commandes
sont rattachés à la cliente en mode \texttt{Restrict}, si bien qu'une
suppression échouerait bruyamment plutôt que de détruire l'historique en
silence. Un compte ainsi marqué ne peut plus se connecter, ni recevoir
d'e-mail, ni réinitialiser son mot de passe.

\paragraph{Le traitement différencié des employées.} Les écritures de fidélité
saisies par le personnel et les bons honorés au comptoir désignent l'employée
concernée, mais en mode \texttt{SetNull} et non \texttt{Cascade}. Le départ
d'une employée ne doit pas effacer les points des clientes qu'elle a servies ni
la trace de ce qui leur a été remis : l'écriture demeure, elle ne désigne
simplement plus personne.

\paragraph{Le bon de récompense comme entité de première classe.} C'est la
notion qui manquait à la première version du programme de fidélité. Sans elle,
réclamer une récompense offerte la faisait disparaître des deux côtés à la fois
--- l'application filtrait sur la date de réclamation, la fiche du comptoir
aussi --- et un échange de points ne laissait qu'une ligne d'historique que
personne ne listait. Le bon réunit les deux circuits de récompense en un objet
unique : une seule chose à afficher pour la cliente, une seule liste à traiter
au comptoir. Son unicité par origine --- un déblocage de palier, ou un
mouvement de points --- constitue le véritable garde-fou d'idempotence : une
réclamation rejouée ne peut pas émettre un second bon.

\paragraph{La plage d'ouverture plutôt que le jour ouvré.} Les horaires sont
modélisés comme des \emph{plages}, et non comme des jours dotés d'une heure
d'ouverture et d'une heure de fermeture. Un même jour en porte donc autant que
nécessaire --- 9\,h--13\,h puis 14\,h--18\,h --- et un jour sans aucune plage
est fermé. Ce modèle exprime naturellement la pause déjeuner, qu'un booléen
« fermé » assorti de deux heures ne savait représenter qu'au prix d'une plage
trouée. Les fermetures exceptionnelles, elles, sont stockées comme des dates de
calendrier et non comme des instants : une fermeture est un jour du calendrier
local du centre, et stocker un instant obligerait à choisir une heure
arbitraire.

\section*{Conclusion}
\addcontentsline{toc}{section}{Conclusion}

Ce chapitre a présenté l'analyse et la conception du système : les besoins
fonctionnels et non fonctionnels, les acteurs, les cas d'utilisation et leur
description, les séquences des processus les plus délicats, et la structure des
données. Les décisions de modélisation y ont été justifiées plutôt que
seulement énoncées, car plusieurs d'entre elles --- le figement des montants,
l'anonymisation, le bon de récompense --- constituent la réponse à des
difficultés qui ne se voyaient qu'à l'usage. Le chapitre suivant expose les
choix techniques qui permettent de réaliser cette conception.
